\input{preamble}

% OK, start here.
%
\begin{document}

\title{Schémas de Picard des courbes}


\maketitle

\phantomsection
\label{section-phantom}

\tableofcontents


\section{Introduction}
\label{section-introduction}

\noindent
Dans ce chapitre, nous nous bornons à construire le schéma de Picard
d'une courbe projective non singulière sur un corps algébriquement clos.
On trouvera dans \cite{Kleiman-Picard} un exposé plus approfondi ainsi que
des indications historiques.

\medskip\noindent
Nous étudierons plus loin, dans le projet Stacks, les foncteurs de Hilbert et de Quot
dans un cadre beaucoup plus général.


\section{Schéma de Hilbert des points}
\label{section-hilbert-scheme-points}

\noindent
Soit $X \to S$ un morphisme de schémas. Soit $d \geq 0$ un entier.
Pour tout schéma $T$ sur $S$, posons
$$
\Hilbfunctor^d_{X/S}(T) =
\left\{
\begin{matrix}
Z \subset X_T\text{ sous-schéma fermé tel que }\\
Z \to T\text{ soit fini localement libre de degré }d
\end{matrix}
\right\}
$$
Si $T' \to T$ est un morphisme de schémas sur $S$ et si
$Z \in \Hilbfunctor^d_{X/S}(T)$, le changement de base
$Z_{T'} \subset X_{T'}$ est un élément de $\Hilbfunctor^d_{X/S}(T')$.
On obtient ainsi un foncteur
$$
\Hilbfunctor^d_{X/S} :
(\Sch/S)^{opp} \longrightarrow \textit{Ensembles},\quad
T \longrightarrow \Hilbfunctor^d_{X/S}(T)
$$
En général, $\Hilbfunctor^d_{X/S}$ est un espace algébrique
(insérer ici une référence ultérieure). Dans cette section, nous
montrerons que $\Hilbfunctor^d_{X/S}$ est représentable
par un schéma si tout ensemble fini de points d'une fibre de
$X \to S$ est contenu dans un ouvert affine.
Lorsque $\Hilbfunctor^d_{X/S}$ est représentable par un schéma, nous
notons souvent ce schéma $\underline{\Hilbfunctor}^d_{X/S}$.

\begin{lemma}
\label{lemma-hilb-d-sheaf}
Soit $X \to S$ un morphisme de schémas. Le foncteur $\Hilbfunctor^d_{X/S}$
est un faisceau pour la topologie fpqc
(Topologies, définition \ref{topologies-definition-sheaf-property-fpqc}).
\end{lemma}

\begin{proof}
Soit $\{T_i \to T\}_{i \in I}$ un recouvrement fpqc de schémas sur $S$.
Posons $X_i = X_{T_i} = X \times_S T_i$.
Notons que $\{X_i \to X_T\}_{i \in I}$ est un recouvrement fpqc de
$X_T$ (Topologies, lemme \ref{topologies-lemma-fpqc})
et que $X_{T_i \times_T T_{i'}} = X_i \times_{X_T} X_{i'}$.
Soit $Z_i \in \Hilbfunctor^d_{X/S}(T_i)$ une famille
d'éléments telle que $Z_i$ et $Z_{i'}$ aient la même image dans
$\Hilbfunctor^d_{X/S}(T_i \times_T T_{i'})$. Par descente effective
des immersions fermées (Descente, lemme \ref{descent-lemma-closed-immersion}),
il existe une immersion fermée $Z \to X_T$ dont le changement de base par
$X_i \to X_T$ est $Z_i \to X_i$. Le morphisme $Z \to T$
a alors pour changement de base à $T_i$ le morphisme
$Z_i \to T_i$. Par conséquent, $Z \to T$ est fini localement libre de degré $d$
d'après Descente, lemme \ref{descent-lemma-descending-property-finite-locally-free}.
\end{proof}

\begin{lemma}
\label{lemma-hilb-d-limit-preserving}
Soit $X \to S$ un morphisme de schémas. Si $X \to S$ est
de présentation finie, le foncteur $\Hilbfunctor^d_{X/S}$
est compatible aux limites projectives (Limites, remarque \ref{limits-remark-limit-preserving}).
\end{lemma}

\begin{proof}
Soit $T = \lim T_i$ une limite de schémas affines sur $S$. Il s'agit de montrer
que $\Hilbfunctor^d_{X/S}(T) = \colim \Hilbfunctor^d_{X/S}(T_i)$.
Si $Z \to X_T$ est un élément de $\Hilbfunctor^d_{X/S}(T)$,
le morphisme $Z \to T$ est de présentation finie. Ainsi, d'après
Limites, lemme \ref{limits-lemma-descend-finite-presentation},
il existe un indice $i$, un schéma $Z_i$ de présentation finie sur $T_i$,
et un morphisme $Z_i \to X_{T_i}$ sur $T_i$ dont le changement de base à $T$
est $Z \to X_T$. Appliquons Limites, lemme
\ref{limits-lemma-descend-closed-immersion-finite-presentation}
pour voir que l'on peut supposer que $Z_i \to X_{T_i}$ est une immersion fermée
quitte à augmenter $i$.
Appliquons Limites, lemme \ref{limits-lemma-descend-finite-locally-free}
pour voir que $Z_i \to T_i$ est fini localement libre de degré $d$
quitte à augmenter encore $i$.
Alors $Z_i \in \Hilbfunctor^d_{X/S}(T_i)$, comme voulu.
\end{proof}

\noindent
Soit $S$ un schéma. Soit $i : X \to Y$ une immersion fermée de schémas
sur $S$. On dispose alors d'une transformation de foncteurs
$$
\Hilbfunctor^d_{X/S} \longrightarrow \Hilbfunctor^d_{Y/S}
$$
qui associe à un élément $Z \in \Hilbfunctor^d_{X/S}(T)$
l'élément $i_T(Z) \subset Y_T$ de $\Hilbfunctor^d_{Y/S}$. Ici, $i_T : X_T \to Y_T$
est le changement de base de $i$.

\begin{lemma}
\label{lemma-hilb-d-of-closed}
Soit $S$ un schéma. Soit $i : X \to Y$ une immersion fermée de schémas.
Si $\Hilbfunctor^d_{Y/S}$ est représentable par un schéma, il en est de même de
$\Hilbfunctor^d_{X/S}$, et le morphisme de schémas correspondant
$\underline{\Hilbfunctor}^d_{X/S} \to \underline{\Hilbfunctor}^d_{Y/S}$
est une immersion fermée.
\end{lemma}

\begin{proof}
Soit $T$ un schéma sur $S$ et soit $Z \in \Hilbfunctor^d_{Y/S}(T)$.
Montrons qu'il existe un sous-schéma fermé $T_X \subset T$ tel
qu'un morphisme de schémas $T' \to T$ se factorise par $T_X$ si
et seulement si $Z_{T'} \to Y_{T'}$ se factorise par $X_{T'}$.
En appliquant ceci à un schéma $T_{univ}$ représentant $\Hilbfunctor^d_{Y/S}$ et à
l'objet universel\footnote{Voir
Catégories, section \ref{categories-section-opposite}}
$Z_{univ} \in \Hilbfunctor^d_{Y/S}(T_{univ})$,
on obtient un sous-schéma fermé $T_{univ, X} \subset T_{univ}$ tel que
$Z_{univ, X} = Z_{univ} \times_{T_{univ}} T_{univ, X}$
soit un sous-schéma fermé de $X \times_S T_{univ, X}$, donc
définisse un élément de $\Hilbfunctor^d_{X/S}(T_{univ, X})$.
Un argument formel montre alors que $T_{univ, X}$ est un schéma
représentant $\Hilbfunctor^d_{X/S}$, d'objet universel $Z_{univ, X}$.

\medskip\noindent
Démontrons l'assertion. Considérons $Z' = X_T \times_{Y_T} Z$. Étant donné $T' \to T$,
on voit que $Z_{T'} \to Y_{T'}$ se factorise par $X_{T'}$ si et
seulement si $Z'_{T'} \to Z_{T'}$ est un isomorphisme. L'assertion résulte donc
du résultat très général
Compléments sur la platitude, lemme \ref{flat-lemma-Weil-restriction-closed-subschemes}.
Dans ce cas particulier, on peut toutefois la démontrer directement comme
suit : on se ramène d'abord au cas $T = \Spec(A)$ et $Z = \Spec(B)$.
En restreignant encore $T$, on peut supposer qu'il existe un isomorphisme
$\varphi : B \to A^{\oplus d}$ de $A$-modules. Alors $Z' = \Spec(B/J)$
pour un idéal $J \subset B$. Soit $g_\beta \in J$ une famille de
générateurs, et écrivons $\varphi(g_\beta) = (g_\beta^1, \ldots, g_\beta^d)$.
Il est alors clair que $T_X$ est donné par $\Spec(A/(g_\beta^j))$.
\end{proof}

\begin{lemma}
\label{lemma-hilb-d-separated}
Soit $X \to S$ un morphisme de schémas. Si $X \to S$ est séparé et si
$\Hilbfunctor^d_{X/S}$ est représentable,
alors $\underline{\Hilbfunctor}^d_{X/S} \to S$ est séparé.
\end{lemma}

\begin{proof}
Dans cette démonstration, tous les produits sans indice sont pris sur $S$.
Posons $H = \underline{\Hilbfunctor}^d_{X/S}$ et soit
$Z \in \Hilbfunctor^d_{X/S}(H)$ l'objet universel.
Considérons les deux objets $Z_1, Z_2 \in \Hilbfunctor^d_{X/S}(H \times H)$
obtenus en prenant l'image réciproque de $Z$ par les deux projections $H \times H \to H$.
Alors $Z_1 = Z \times H \subset X_{H \times H}$ et $Z_2 = H \times Z
\subset X_{H \times H}$. Puisque $H$ représente le foncteur
$\Hilbfunctor^d_{X/S}$, le morphisme diagonal $\Delta : H \to H \times H$
a la propriété universelle suivante : un morphisme de schémas
$T \to H \times H$ se factorise par $\Delta$ si et seulement si
$Z_{1, T} = Z_{2, T}$ en tant qu'éléments de $\Hilbfunctor^d_{X/S}(T)$.
Posons $Z = Z_1 \times_{X_{H \times H}} Z_2$. On voit alors que
$T \to H \times H$ se factorise par $\Delta$ si et seulement si
les morphismes $Z_T \to Z_{1, T}$ et $Z_T \to Z_{2, T}$ sont des
isomorphismes. Le résultat très général
Compléments sur la platitude, lemme \ref{flat-lemma-Weil-restriction-closed-subschemes}
montre que $\Delta$ est une immersion fermée. Dans la démonstration du
lemme \ref{lemma-hilb-d-of-closed},
le lecteur trouvera une autre démonstration, plus simple, du résultat nécessaire
dans notre cas particulier.
\end{proof}

\begin{lemma}
\label{lemma-hilb-d-An}
Soit $X \to S$ un morphisme de schémas affines. Soit $d \geq 0$. Alors
$\Hilbfunctor^d_{X/S}$ est représentable.
\end{lemma}

\begin{proof}
Écrivons $S = \Spec(R)$. On peut choisir une immersion fermée de $X$
dans le spectre de $R[x_i; i \in I]$ pour un ensemble $I$ (de cardinal
assez grand. Ainsi, d'après le lemme \ref{lemma-hilb-d-of-closed},
on peut supposer que $X = \Spec(A)$, où $A = R[x_i; i \in I]$.
Nous utiliserons Schémas, lemme \ref{schemes-lemma-glue-functors} pour démontrer le
lemme dans ce cas.

\medskip\noindent
La condition (1) du lemme résulte du lemme \ref{lemma-hilb-d-sheaf}.

\medskip\noindent
Pour toute partie $W \subset A$ de cardinal $d$, nous allons
construire un sous-foncteur $F_W$ de $\Hilbfunctor^d_{X/S}$.
(Il suffirait de considérer le cas où $W$ est formé de
monômes en les $x_i$, mais cela ne nous sera pas nécessaire.)
Nous dirons que $Z \in \Hilbfunctor^d_{X/S}(T)$ appartient à $F_W(T)$
si et seulement si l'application $\mathcal{O}_T$-linéaire
$$
\bigoplus\nolimits_{f \in W} \mathcal{O}_T
\longrightarrow
(Z \to T)_*\mathcal{O}_Z,\quad
(g_f) \longmapsto \sum g_f f|_Z
$$
est surjective (ou, ce qui revient au même, est un isomorphisme). Pour $f \in A$
et $Z \in \Hilbfunctor^d_{X/S}(T)$, on note ici $f|_Z$ l'image réciproque de $f$
par le morphisme $Z \to X_T \to X$.

\medskip\noindent
Ouverture, c'est-à-dire condition (2)(b) du lemme. Elle résulte de
Algèbre, lemme \ref{algebra-lemma-cokernel-flat}.

\medskip\noindent
Recouvrement, c'est-à-dire condition (2)(c) du lemme. Puisque
$$
A \otimes_R \mathcal{O}_T =
(X_T \to T)_*\mathcal{O}_{X_T} \to (Z \to T)_*\mathcal{O}_Z
$$
est surjective et que $(Z \to T)_*\mathcal{O}_Z$ est localement libre
de rang fini $d$, on peut, pour tout point $t \in T$, trouver une partie finie
$W \subset A$ de cardinal $d$ dont les images forment une base
de l'espace vectoriel de dimension $d$ sur $\kappa(t)$
$((Z \to T)_*\mathcal{O}_Z)_t \otimes_{\mathcal{O}_{T, t}} \kappa(t)$.
Le lemme de Nakayama donne un voisinage ouvert $V \subset T$
de $t$ tel que $Z_V \in F_W(V)$.

\medskip\noindent
Représentabilité, c'est-à-dire condition (2)(a) du lemme. Soit $W \subset A$
de cardinal $d$. Nous affirmons que $F_W$ est représentable par un schéma
affine sur $R$. Nous allons construire ce schéma affine, mais invitons
le lecteur à en chercher lui-même la construction. Numérotons
$f_1, \ldots, f_d$ les éléments de $W$. Nous construirons un élément
universel $Z_{univ} = \Spec(B_{univ})$ de $F_W$ sur $T_{univ} = \Spec(R_{univ})$
comme suit. Posons d'abord
$$
R_{univ} = R[c_{kl}^m, b^l, b_i^l]/\mathfrak a_{univ}
$$
pour un idéal $\mathfrak a_{univ}$ que nous décrirons ci-dessous.
Les indices $k, l, m, i$ parcourent tous ici $\{1, \ldots, d\}$.
Posons ensuite
$$
B_{univ} = R_{univ}e_1 \oplus R_{univ}e_2 \oplus \ldots \oplus R_{univ}e_d
$$
en tant que $R_{univ}$-module, muni de la structure d'algèbre donnée par
$$
e_ke_l = \sum c_{kl}^m e_m
$$
L'immersion fermée $Z_{univ} \to X_{T_{univ}}$ est donnée
par l'homomorphisme de $R_{univ}$-algèbres
$$
\Psi : A \otimes_R R_{univ} \longrightarrow B_{univ}
$$
envoyant $1 \otimes 1$ sur $\sum b^le_l$ et $x_i$ sur $\sum b_i^le_l$
(voir ci-dessous pour son existence). Quant à l'idéal $\mathfrak a_{univ}$,
il doit satisfaire aux prescriptions suivantes :
\begin{enumerate}
\item la multiplication dans $B_{univ}$ est commutative, c'est-à-dire que
l'on doit avoir $c_{lk}^m - c_{kl}^m \in \mathfrak a_{univ}$,
\item la multiplication dans $B_{univ}$ est associative, c'est-à-dire que
l'on doit avoir $c_{lk}^m c_{m n}^p - c_{lq}^p c_{kn}^q \in \mathfrak a_{univ}$,
\item $\sum b^le_l$ est un élément unité $1$ de $B_{univ}$,\linebreak
autrement dit, on doit avoir $(\sum b^le_l)e_k = e_k$ pour tout $k$,
ce qui signifie que $\sum b^lc_{lk}^m - \delta_{km} \in \mathfrak a_{univ}$
(symbole de Kronecker).
\end{enumerate}
Puisque nous avons assuré que $B_{univ}$ est une
$R_{univ}$-algèbre commutative dont l'unité $1$ est $\sum b^le_l$, l'homomorphisme $\Psi$
existe. La dernière condition est la suivante :
\begin{enumerate}
\item[(5)] $f_l$ doit avoir pour image $e_l$ dans $B_{univ}$. Écrivons
$\Psi(f_l) - e_l \equiv \sum h_l^me_m$ pour des
$h_l^m \in R[c_{kl}^m, b^l, b_i^l]$ ;
on doit alors avoir $h_l^m \in \mathfrak a_{univ}$.
\end{enumerate}
En prenant pour $\mathfrak a_{univ} \subset R[c_{kl}^m, b^l, b_i^l]$ l'idéal engendré
par les éléments énumérés dans (1) -- (5),
il est clair que $F_W$ est représenté par $\Spec(R_{univ})$.
\end{proof}

\begin{proposition}
\label{proposition-hilb-d-representable}
Soit $X \to S$ un morphisme de schémas. Soit $d \geq 0$. Supposons
que pour tout $(s, x_1, \ldots, x_d)$, où $s \in S$ et
$x_1, \ldots, x_d \in X_s$, il existe un ouvert affine $U \subset X$
tel que $x_1, \ldots, x_d \in U$. Alors $\Hilbfunctor^d_{X/S}$ est
représentable par un schéma.
\end{proposition}

\begin{proof}
En utilisant soit le recollement relatif (Constructions, section
\ref{constructions-section-relative-glueing}), soit
le point de vue fonctoriel
(Schémas, lemme \ref{schemes-lemma-glue-functors}),
on se ramène au cas où $S$ est affine. Nous omettons les détails.

\medskip\noindent
Supposons $S$ affine. Pour un ouvert affine $U \subset X$, notons
$F_U \subset \Hilbfunctor^d_{X/S}$ le sous-foncteur tel que,
pour tout schéma $T/S$, un élément $Z \in \Hilbfunctor^d_{X/S}(T)$
appartienne à $F_U(T)$ si et seulement si $Z \subset U_T$. Nous utiliserons
Schémas, lemme \ref{schemes-lemma-glue-functors}
et les sous-foncteurs $F_U$ pour conclure.

\medskip\noindent
La condition (1) est le lemme \ref{lemma-hilb-d-sheaf}.

\medskip\noindent
La condition (2)(a) résulte de ce que $F_U = \Hilbfunctor^d_{U/S}$
et de sa représentabilité, donnée par le lemme \ref{lemma-hilb-d-An}.
En effet, si $Z \in F_U(T)$, on peut regarder $Z$ comme un sous-schéma fermé
de $U_T$, fini localement libre de degré $d$ sur $T$, d'où
$Z \in \Hilbfunctor^d_{U/S}(T)$. Réciproquement, si $Z \in \Hilbfunctor^d_{U/S}(T)$,
alors $Z \to U_T \to X_T$ est une immersion fermée\footnote{C'est clair
si $X \to S$ est séparé, car dans ce cas Morphismes, lemme
\ref{morphisms-lemma-image-proper-scheme-closed}
montre que l'immersion $\varphi : Z \to X_T$ a une image fermée,
et est donc une immersion fermée d'après
Schémas, lemme \ref{schemes-lemma-immersion-when-closed}. Nous suggérons au
lecteur de passer la suite de cette note, car nous ne connaissons aucun exemple
où les hypothèses sur $X \to S$ soient satisfaites sans que $X \to S$ soit séparé.
Dans le cas général, soit $x \in X_T$ un point de l'adhérence de
$\varphi(Z)$. Il faut montrer que $x \in \varphi(Z)$. Soit $t \in T$ l'image
de $x$. L'hypothèse sur $X \to S$ permet de choisir un ouvert affine
$W \subset X_T$ contenant $x$ et $\varphi(Z_t)$. Alors $\varphi^{-1}(W)$
est un ouvert contenant toute la fibre $Z_t$ et, puisque $Z \to T$ est fermé,
on peut, en remplaçant $T$ par un voisinage ouvert de $t$, supposer que
$Z = \varphi^{-1}(W)$. Alors $\varphi(Z) \subset W$ est fermé d'après le
cas séparé (car $W \to T$ est séparé), d'où $x \in \varphi(Z)$.}
et l'on peut regarder $Z$ comme un élément de $F_U(T)$.

\medskip\noindent
Soit $Z \in \Hilbfunctor^d_{X/S}(T)$ pour un schéma $T$ sur $S$. Posons
$$
B = (Z \to T)\left((Z \to X_T \to X)^{-1}(X \setminus U)\right)
$$
C'est une partie fermée de $T$, et il est clair que sur l'ouvert
$T_{Z, U} = T \setminus B$, la restriction $Z_{t'}$ a son image dans $U_{T'}$.
D'autre part, pour tout $b \in B$, la fibre $Z_b$ n'a pas son image
dans $U$. Ainsi, étant donné un morphisme $T' \to T$, on
a $Z_{T'} \in F_U(T')$ $\Leftrightarrow$ $T' \to T$ se factorise par
l'ouvert $T_{Z, U}$. Cela démontre la condition (2)(b).

\medskip\noindent
La condition (2)(c) résulte de notre hypothèse sur $X/S$. Il suffit de
montrer ceci : si $T$ est le spectre d'un corps
et si $Z \subset X_T$ est un sous-schéma fermé, fini plat de degré
$d$ sur $T$, alors $Z \to X_T \to X$ se factorise par un ouvert affine
$U$ de $X$. C'est clair, car $Z$ a au plus $d$ points,
dont les images sont toutes dans la fibre de $X$ au-dessus du point image
de $T \to S$.
\end{proof}

\begin{remark}
\label{remark-when-proposition-applies}
Soit $f : X \to S$ un morphisme de schémas. L'hypothèse de la
proposition \ref{proposition-hilb-d-representable}, et par
suite sa conclusion, sont satisfaites dans chacun des cas suivants :
\begin{enumerate}
\item $X$ est quasi-affine,
\item $f$ est quasi-affine,
\item $f$ est quasi-projectif,
\item $f$ est localement projectif,
\item il existe un faisceau inversible ample sur $X$,
\item il existe un faisceau inversible $f$-ample sur $X$, et
\item il existe un faisceau inversible $f$-très ample sur $X$.
\end{enumerate}
En effet, dans chacun de ces cas, tout ensemble fini de points d'une
fibre $X_s$ est contenu dans un ouvert quasi-compact $U$ de $X$
qui est muni d'un faisceau inversible ample, est isomorphe
à un ouvert d'un schéma affine, ou est isomorphe à un ouvert
du $\text{Proj}$ d'un anneau gradué (dans chaque cas, cela résulte
directement des définitions). L'existence d'ouverts affines convenables
résulte donc de
Propriétés, lemme \ref{properties-lemma-ample-finite-set-in-affine}.
\end{remark}




\section{Espaces de modules de diviseurs sur les courbes lisses}
\label{section-divisors}

\noindent
Pour un morphisme lisse $X \to S$ de dimension relative $1$, le foncteur
$\Hilbfunctor^d_{X/S}$ paramètre les diviseurs de Cartier effectifs relatifs
au sens de
Diviseurs, section \ref{divisors-section-effective-Cartier-morphisms}.

\begin{lemma}
\label{lemma-divisors-on-curves}
Soit $X \to S$ un morphisme lisse de schémas de dimension relative $1$.
Soit $D \subset X$ un sous-schéma fermé. Considérons les conditions suivantes :
\begin{enumerate}
\item $D \to S$ est fini localement libre,
\item $D$ est un diviseur de Cartier effectif relatif sur $X/S$,
\item $D \to S$ est localement quasi-fini, plat et
localement de présentation finie, et
\item $D \to S$ est localement quasi-fini et plat.
\end{enumerate}
On a toujours les implications
$$
(1) \Rightarrow (2) \Leftrightarrow (3) \Rightarrow (4)
$$
Si $S$ est localement noethérien, la dernière flèche est une équivalence.
Si $X \to S$ est propre (et $S$ quelconque), la première flèche est
une équivalence.
\end{lemma}

\begin{proof}
Équivalence de (2) et (3). Elle résulte de
Diviseurs, lemme \ref{divisors-lemma-fibre-Cartier},
à condition de montrer l'équivalence de (2) et (3) lorsque
$S$ est le spectre d'un corps $k$. Soit $x \in X$ un point fermé.
Puisque $X$ est lisse de dimension relative $1$ sur $k$, on voit que
$\mathcal{O}_{X, x}$ est un anneau local régulier de dimension $1$
(voir Variétés, lemme \ref{varieties-lemma-smooth-regular}).
Ainsi, $\mathcal{O}_{X, x}$ est un anneau de valuation discrète
(Algèbre, lemme \ref{algebra-lemma-characterize-dvr}),
donc un anneau principal. Il en résulte que tout faisceau d'idéaux
$\mathcal{I} \subset \mathcal{O}_X$ non nul en tous
les points génériques de $X$ est inversible
(Diviseurs, lemme \ref{divisors-lemma-effective-Cartier-in-points}).
Autrement dit, tout sous-schéma fermé de $X$ ne contenant
aucun point générique est un diviseur de Cartier effectif.
Il s'ensuit que (2) et (3) sont équivalentes.

\medskip\noindent
Si $S$ est noethérien, tout morphisme localement quasi-fini
$D \to S$ est localement de présentation finie (Morphismes, lemme
\ref{morphisms-lemma-noetherian-finite-type-finite-presentation}),
d'où l'équivalence de (3) et (4).

\medskip\noindent
Si $X \to S$ est propre (et $S$ quelconque), alors $D \to S$ est
également propre. Un morphisme propre localement quasi-fini étant fini
(Compléments sur les morphismes, lemme \ref{more-morphisms-lemma-characterize-finite})
et un morphisme fini, plat et de présentation finie étant fini localement libre
(Morphismes, lemme \ref{morphisms-lemma-finite-flat}), on voit que
(1) et (2) sont équivalentes.
\end{proof}

\begin{lemma}
\label{lemma-sum-divisors-on-curves}
Soit $X \to S$ un morphisme lisse de schémas de dimension relative $1$.
Soient $D_1, D_2 \subset X$ des sous-schémas fermés finis localement libres de
degrés $d_1$, $d_2$ sur $S$. Alors $D_1 + D_2$ est fini localement libre
de degré $d_1 + d_2$ sur $S$.
\end{lemma}

\begin{proof}
D'après le lemme \ref{lemma-divisors-on-curves}, les sous-schémas $D_1$
et $D_2$ sont des diviseurs de Cartier effectifs relatifs sur $X/S$.
Ainsi, $D = D_1 + D_2$ est un diviseur de Cartier effectif relatif
sur $X/S$ d'après
Diviseurs, lemme \ref{divisors-lemma-sum-relative-effective-Cartier-divisor}.
Par conséquent, $D \to S$ est localement quasi-fini, plat et
localement de présentation finie d'après le
lemme \ref{lemma-divisors-on-curves}.
En appliquant
Morphismes, lemme \ref{morphisms-lemma-image-universally-closed-separated}
au morphisme entier surjectif $D_1 \amalg D_2 \to D$,
on voit que $D \to S$ est séparé. Alors
Morphismes, lemme \ref{morphisms-lemma-image-proper-is-proper}
montre que $D \to S$ est propre.
Il en résulte que $D \to S$ est fini
(Compléments sur les morphismes, lemme \ref{more-morphisms-lemma-characterize-finite}),
puis que $D \to S$ est fini localement libre
(Morphismes, lemme \ref{morphisms-lemma-finite-flat}).
Il suffit donc de montrer que le degré de $D \to S$ est $d_1 + d_2$.
Pour cela, on peut effectuer un changement de base à une fibre de $X \to S$, donc
supposer que $S = \Spec(k)$ pour un corps $k$.
Dans ce cas, il existe un ensemble fini de points fermés
$x_1, \ldots, x_n \in X$ tel que les supports de $D_1$ et $D_2$
soient contenus dans $\{x_1, \ldots, x_n\}$.
Il existe en fait des éléments non diviseurs de zéro $f_{i, j} \in \mathcal{O}_{X, x_i}$
tels que
$$
D_1 = \coprod \Spec(\mathcal{O}_{X, x_i}/(f_{i, 1}))
\quad\text{et}\quad
D_2 = \coprod \Spec(\mathcal{O}_{X, x_i}/(f_{i, 2}))
$$
On voit alors que
$$
D = \coprod \Spec(\mathcal{O}_{X, x_i}/(f_{i, 1}f_{i, 2}))
$$
On en déduit aisément que $D$ est de degré $d_1 + d_2$
sur $k$ (au besoin, utiliser Algèbre, lemme \ref{algebra-lemma-ord-additive}).
\end{proof}

\begin{lemma}
\label{lemma-difference-divisors-on-curves}
Soit $X \to S$ un morphisme lisse de schémas de dimension relative $1$.
Soient $D_1, D_2 \subset X$ des sous-schémas fermés finis localement libres de
degrés $d_1$, $d_2$ sur $S$. Si $D_1 \subset D_2$ (comme sous-schémas fermés),
il existe un sous-schéma fermé $D \subset X$ fini localement libre de
degré $d_2 - d_1$ sur $S$ tel que $D_2 = D_1 + D$.
\end{lemma}

\begin{proof}
Cette démonstration est presque identique à celle du
lemme \ref{lemma-sum-divisors-on-curves}.
D'après le lemme \ref{lemma-divisors-on-curves}, les sous-schémas $D_1$
et $D_2$ sont des diviseurs de Cartier effectifs relatifs sur $X/S$.
D'après Diviseurs, lemme
\ref{divisors-lemma-difference-relative-effective-Cartier-divisor},
il existe un diviseur de Cartier effectif relatif $D \subset X$
tel que $D_2 = D_1 + D$. Ainsi, $D \to S$ est localement quasi-fini, plat et
localement de présentation finie d'après le
lemme \ref{lemma-divisors-on-curves}.
Puisque $D$ est un sous-schéma fermé de $D_2$, on voit que
$D \to S$ est fini. Il s'ensuit que $D \to S$ est fini localement libre
(Morphismes, lemme \ref{morphisms-lemma-finite-flat}).
Il suffit donc de montrer que le degré de $D \to S$ est $d_2 - d_1$.
Cela résulte du lemme \ref{lemma-sum-divisors-on-curves}.
\end{proof}

\noindent
Soit $X \to S$ un morphisme lisse de schémas de dimension relative $1$.
D'après le lemme \ref{lemma-divisors-on-curves}, pour un schéma $T$ sur $S$ et
$D \in \Hilbfunctor^d_{X/S}(T)$, on peut regarder $D$ comme un diviseur de Cartier
effectif relatif sur $X_T/T$ tel que $D \to T$ soit fini
localement libre de degré $d$. Par conséquent, le
lemme \ref{lemma-sum-divisors-on-curves} donne une transformation
de foncteurs
$$
\Hilbfunctor^{d_1}_{X/S} \times \Hilbfunctor^{d_2}_{X/S}
\longrightarrow
\Hilbfunctor^{d_1 + d_2}_{X/S},\quad
(D_1, D_2) \longmapsto D_1 + D_2
$$
Si $\Hilbfunctor^d_{X/S}$ est représentable pour tout degré $d$, cette
transformation de foncteurs correspond à un morphisme de schémas
$$
\underline{\Hilbfunctor}^{d_1}_{X/S}
\times_S
\underline{\Hilbfunctor}^{d_2}_{X/S}
\longrightarrow
\underline{\Hilbfunctor}^{d_1 + d_2}_{X/S}
$$
sur $S$. Notons que $\underline{\Hilbfunctor}^0_{X/S} = S$ et
$\underline{\Hilbfunctor}^1_{X/S} = X$.
Un cas particulier du morphisme précédent est le morphisme
$$
\underline{\Hilbfunctor}^d_{X/S} \times_S X
\longrightarrow
\underline{\Hilbfunctor}^{d + 1}_{X/S},\quad
(D, x) \longmapsto D + x
$$

\begin{lemma}
\label{lemma-universal-object}
Soit $X \to S$ un morphisme lisse de schémas de dimension relative $1$
tel que les foncteurs $\Hilbfunctor^d_{X/S}$ soient représentables. Le morphisme
$\underline{\Hilbfunctor}^d_{X/S} \times_S X \to
\underline{\Hilbfunctor}^{d + 1}_{X/S}$
est fini localement libre de degré $d + 1$.
\end{lemma}

\begin{proof}
Soit $D_{univ} \subset X \times_S \underline{\Hilbfunctor}^{d + 1}_{X/S}$
l'objet universel. On a un diagramme commutatif
$$
\xymatrix{
\underline{\Hilbfunctor}^d_{X/S} \times_S X \ar[rr] \ar[rd] & &
D_{univ} \ar[ld] \ar@{^{(}->}[r] &
\underline{\Hilbfunctor}^{d + 1}_{X/S} \times_S X \\
& \underline{\Hilbfunctor}^{d + 1}_{X/S}
}
$$
où la flèche horizontale supérieure envoie $(D', x)$ sur $(D' + x, x)$.
Nous affirmons que ce morphisme est un isomorphisme,
ce qui démontre le lemme. En effet, étant donné un schéma $T$ sur $S$,
un $T$-point $\xi$ de $D_{univ}$ est donné par un couple $\xi = (D, x)$,
où $D \subset X_T$ est un sous-schéma fermé fini localement libre
de degré $d + 1$ sur $T$ et $x : T \to X$ est un morphisme dont
le graphe $x : T \to X_T$ se factorise par $D$. D'après le
lemme \ref{lemma-difference-divisors-on-curves},
on peut alors écrire $D = D' + x$ pour un $D' \subset X_T$ fini localement
libre de degré $d$ sur $T$. L'application qui associe à $\xi = (D, x)$ le couple
$(D', x)$ est l'inverse cherché.
\end{proof}

\begin{lemma}
\label{lemma-hilb-d-smooth}
Soit $X \to S$ un morphisme lisse de schémas de dimension relative $1$
tel que les foncteurs $\Hilbfunctor^d_{X/S}$ soient représentables. Les
schémas $\underline{\Hilbfunctor}^d_{X/S}$ sont lisses sur $S$ de
dimension relative $d$.
\end{lemma}

\begin{proof}
On a $\underline{\Hilbfunctor}^0_{X/S} = S$ et
$\underline{\Hilbfunctor}^1_{X/S} = X$ ; le résultat est donc vrai pour $d = 0, 1$.
Supposons le résultat établi pour $d$ ; alors
$\underline{\Hilbfunctor}^d_{X/S} \times_S X$ est lisse sur $S$
(Morphismes, lemme \ref{morphisms-lemma-base-change-smooth} et
\ref{morphisms-lemma-composition-smooth}). Puisque
$\underline{\Hilbfunctor}^d_{X/S} \times_S X \to
\underline{\Hilbfunctor}^{d + 1}_{X/S}$
est fini localement libre de degré $d + 1$ d'après le
lemme \ref{lemma-universal-object},
le résultat découle de
Descente, lemme \ref{descent-lemma-smooth-permanence}.
Nous omettons la vérification de la valeur annoncée de la dimension
relative (on peut la faire en examinant les fibres, ou en
suivant les dimensions dans le raisonnement précédent).
\end{proof}

\noindent
Rassemblons les résultats obtenus jusqu'ici dans le cas
d'une courbe propre et lisse sur un corps.

\begin{proposition}
\label{proposition-hilb-d}
Soit $X$ une courbe propre, lisse et géométriquement irréductible sur un corps $k$.
\begin{enumerate}
\item Les foncteurs $\Hilbfunctor^d_{X/k}$ sont représentables par des variétés
propres et lisses $\underline{\Hilbfunctor}^d_{X/k}$ de dimension
$d$ sur $k$.
\item Pour toute extension de corps $k'/k$, les points $k'$-rationnels
de $\underline{\Hilbfunctor}^d_{X/k}$ sont en bijection $1$-à-$1$
avec les diviseurs de Cartier effectifs de degré $d$ sur $X_{k'}$.
\item Pour $d_1, d_2 \geq 0$, il existe un morphisme
$$
\underline{\Hilbfunctor}^{d_1}_{X/k}
\times_k
\underline{\Hilbfunctor}^{d_2}_{X/k}
\longrightarrow
\underline{\Hilbfunctor}^{d_1 + d_2}_{X/k}
$$
qui est fini localement libre de degré ${d_1 + d_2 \choose d_1}$.
\end{enumerate}
\end{proposition}

\begin{proof}
Les foncteurs $\Hilbfunctor^d_{X/k}$ sont représentables d'après la
proposition \ref{proposition-hilb-d-representable}
(voir aussi la remarque \ref{remark-when-proposition-applies})
et le fait que $X$ est projectif
(Variétés, lemme \ref{varieties-lemma-dim-1-proper-projective}).
Les schémas $\underline{\Hilbfunctor}^d_{X/k}$ sont séparés
sur $k$ d'après le lemme \ref{lemma-hilb-d-separated}.
Les schémas $\underline{\Hilbfunctor}^d_{X/k}$ sont lisses
sur $k$ d'après le lemme \ref{lemma-hilb-d-smooth}.
Partant de $X = \underline{\Hilbfunctor}^1_{X/k}$,
les morphismes du lemme \ref{lemma-universal-object}
donnent par récurrence un morphisme
$$
X^d = X \times_k X \times_k \ldots \times_k X \longrightarrow
\underline{\Hilbfunctor}^d_{X/k},\quad
(x_1, \ldots, x_d) \longrightarrow x_1 + \ldots + x_d
$$
qui est fini localement libre de degré $d!$. Puisque $X$ est
propre sur $k$, il en est de même de $X^d$, donc
$\underline{\Hilbfunctor}^d_{X/k}$ est propre sur $k$ d'après
Morphismes, lemme \ref{morphisms-lemma-image-proper-is-proper}.
Puisque $X$ est géométriquement irréductible sur $k$, le produit
$X^d$ est irréductible
(Variétés, lemme \ref{varieties-lemma-bijection-irreducible-components}),
donc son image est irréductible (et même géométriquement irréductible).
Cela démontre (1). L'assertion (2) résulte des définitions. L'assertion (3) découle
du diagramme commutatif
$$
\xymatrix{
X^{d_1} \times_k X^{d_2} \ar[d] \ar@{=}[r] & X^{d_1 + d_2} \ar[d] \\
\underline{\Hilbfunctor}^{d_1}_{X/k}
\times_k
\underline{\Hilbfunctor}^{d_2}_{X/k}
\ar[r] &
\underline{\Hilbfunctor}^{d_1 + d_2}_{X/k}
}
$$
et de la multiplicativité des degrés des morphismes finis localement libres.
\end{proof}

\begin{remark}
\label{remark-universal-object-hilb-d}
Soit $X$ une courbe propre, lisse et géométriquement irréductible sur un corps $k$,
comme dans la proposition \ref{proposition-hilb-d}. Soit $d \geq 0$. L'objet
fermé universel est un diviseur effectif relatif
$$
D_{univ} \subset \underline{\Hilbfunctor}^{d + 1}_{X/k} \times_k X
$$
sur $\underline{\Hilbfunctor}^{d + 1}_{X/k}$ d'après le
lemme \ref{lemma-divisors-on-curves}.
En fait, $D_{univ}$ est isomorphe, en tant que schéma, à
$\underline{\Hilbfunctor}^d_{X/k} \times_k X$, voir la démonstration du
lemme \ref{lemma-universal-object}.
En particulier, $D_{univ}$ est un diviseur de Cartier effectif et
l'on obtient un module inversible
$\mathcal{O}(D_{univ})$. Si $[D] \in \underline{\Hilbfunctor}^{d + 1}_{X/k}$
désigne le point $k$-rationnel correspondant au diviseur de Cartier
effectif $D \subset X$ de degré $d + 1$, la restriction
de $\mathcal{O}(D_{univ})$ à la fibre $[D] \times X$ est
$\mathcal{O}_X(D)$.
\end{remark}


\section{Le foncteur de Picard}
\label{section-picard-functor}

\noindent
Pour tout schéma $X$, on note $\Pic(X)$ l'ensemble des classes
d'isomorphisme de $\mathcal{O}_X$-modules inversibles.
Voir Modules, définition \ref{modules-definition-pic}.
Étant donné un morphisme $f : X \to Y$ de schémas, l'image réciproque définit
un homomorphisme de groupes $\Pic(Y) \to \Pic(X)$.
L'association
$X \leadsto \Pic(X)$ est un foncteur contravariant de la catégorie
des schémas dans celle des groupes abéliens. Ce foncteur n'est pas
représentable, mais une variante relative de cette
construction est parfois représentable.

\medskip\noindent
Définissons le foncteur de Picard d'un morphisme de schémas $f : X \to S$.
L'idée de la construction consiste à prendre pour ce foncteur le faisceau
$R^1f_*\mathbf{G}_m$, l'image directe supérieure étant calculée pour
la topologie fppf. En explicitant les définitions, on obtient la définition
plus directe suivante.

\begin{definition}
\label{definition-picard-functor}
Soit $\Sch_{fppf}$ un grand site comme dans
Topologies, définition \ref{topologies-definition-big-small-fppf}.
Soit $f : X \to S$ un morphisme de ce site. Le {\it foncteur de Picard}
$\Picardfunctor_{X/S}$ est le faisceau fppf associé au foncteur
$$
(\Sch/S)_{fppf} \longrightarrow \textit{Ensembles},\quad
T \longmapsto \Pic(X_T)
$$
Si ce foncteur est représentable, on note
$\underline{\Picardfunctor}_{X/S}$ un schéma qui le représente.
\end{definition}

\noindent
Une remarque d'usage fréquent est que, si $T \in \Ob((\Sch/S)_{fppf})$, alors
$\Picardfunctor_{X_T/T}$ est la restriction de $\Picardfunctor_{X/S}$ à
$(\Sch/T)_{fppf}$.
Déterminer la valeur de $\Picardfunctor_{X/S}$ sur les schémas
$T$ sur $S$ n'est pas immédiat. Le lemme suivant facilite cette
tâche.

\begin{lemma}
\label{lemma-flat-geometrically-connected-fibres}
Soit $f : X \to S$ comme dans la définition \ref{definition-picard-functor}.
Si $\mathcal{O}_T \to f_{T, *}\mathcal{O}_{X_T}$ est un isomorphisme
pour tout $T \in \Ob((\Sch/S)_{fppf})$, alors
$$
0 \to \Pic(T) \to \Pic(X_T) \to \Picardfunctor_{X/S}(T)
$$
est une suite exacte pour tout $T$.
\end{lemma}

\begin{proof}
On peut remplacer $S$ par $T$ et $X$ par $X_T$, donc supposer $S = T$
pour simplifier les notations. Soit $\mathcal{N}$ un
$\mathcal{O}_S$-module inversible. Si $f^*\mathcal{N} \cong \mathcal{O}_X$, alors
$f_*f^*\mathcal{N} \cong f_*\mathcal{O}_X \cong \mathcal{O}_S$
par hypothèse. Puisque $\mathcal{N}$ est localement trivial, on voit que
l'application canonique $\mathcal{N} \to f_*f^*\mathcal{N}$ est localement
un isomorphisme (car $\mathcal{O}_S \to f_*f^*\mathcal{O}_S$
est un isomorphisme par hypothèse). On en conclut que
$\mathcal{N} \to f_*f^*\mathcal{N} \to \mathcal{O}_S$ est un isomorphisme,
donc que $\mathcal{N}$ est trivial. Cela démontre que la première flèche
est injective.

\medskip\noindent
Soit $\mathcal{L}$ un $\mathcal{O}_X$-module inversible dont la classe appartient
au noyau de $\Pic(X) \to \Picardfunctor_{X/S}(S)$. Il existe alors
un recouvrement fppf $\{S_i \to S\}$ tel que l'image réciproque de $\mathcal{L}$
soit le faisceau inversible trivial sur $X_{S_i}$. Choisissons une section
trivialisante $s_i$. Alors $\text{pr}_0^*s_i$ et $\text{pr}_1^*s_j$ sont deux
sections trivialisantes de $\mathcal{L}$ sur $X_{S_i \times_S S_j}$
et diffèrent donc par multiplication par une unité
$$
f_{ij} \in
\Gamma(X_{S_i \times_S S_j}, \mathcal{O}_{X_{S_i \times_S S_j}}^*) =
\Gamma(S_i \times_S S_j, \mathcal{O}_{S_i \times_S S_j}^*)
$$
(l'égalité résulte de notre hypothèse sur l'image directe des faisceaux structuraux).
Ces éléments satisfont évidemment à la condition de cocycle sur
$S_i \times_S S_j \times_S S_k$, et définissent donc une donnée de descente
de faisceaux inversibles pour le recouvrement fppf $\{S_i \to S\}$.
D'après Descente, proposition \ref{descent-proposition-fpqc-descent-quasi-coherent},
il existe un $\mathcal{O}_S$-module inversible $\mathcal{N}$
muni de trivialisations sur les $S_i$ dont la donnée de descente associée est
$\{f_{ij}\}$. Alors $f^*\mathcal{N} \cong \mathcal{L}$, car le
foncteur des données de descente vers les modules est pleinement fidèle (voir la proposition
citée ci-dessus).
\end{proof}

\begin{lemma}
\label{lemma-flat-geometrically-connected-fibres-with-section}
Soit $f : X \to S$ comme dans la définition \ref{definition-picard-functor}.
Supposons que $f$ admette une section $\sigma$ et que
$\mathcal{O}_T \to f_{T, *}\mathcal{O}_{X_T}$ soit un isomorphisme
pour tout $T \in \Ob((\Sch/S)_{fppf})$. Alors
$$
0 \to \Pic(T) \to \Pic(X_T) \to \Picardfunctor_{X/S}(T) \to 0
$$
est une suite exacte scindée, le scindage étant donné par
$\sigma_T^* : \Pic(X_T) \to \Pic(T)$.
\end{lemma}

\begin{proof}
Notons $K(T) = \Ker(\sigma_T^* : \Pic(X_T) \to \Pic(T))$.
Puisque $\sigma$ est une section de $f$, le groupe $\Pic(X_T)$ est la somme
directe de $\Pic(T)$ et de $K(T)$.
Le lemme \ref{lemma-flat-geometrically-connected-fibres} montre donc que
$K(T) \subset \Picardfunctor_{X/S}(T)$ pour tout $T$. De plus, la construction
montre clairement que $\Picardfunctor_{X/S}$ est le faisceau associé
au préfaisceau $K$. Pour achever la démonstration, il suffit de montrer que
$K$ satisfait à la condition de faisceau pour les recouvrements fppf, ce que nous faisons
dans le paragraphe suivant.

\medskip\noindent
Soit $\{T_i \to T\}$ un recouvrement fppf. Soient $\mathcal{L}_i$ des
éléments de $K(T_i)$ dont les images dans $K(T_i \times_T T_j)$ sont égales
pour tous $i$ et $j$. Choisissons un isomorphisme
$\alpha_i : \mathcal{O}_{T_i} \to \sigma_{T_i}^*\mathcal{L}_i$
pour tout $i$. Choisissons un isomorphisme
$$
\varphi_{ij} :
\mathcal{L}_i|_{X_{T_i \times_T T_j}}
\longrightarrow
\mathcal{L}_j|_{X_{T_i \times_T T_j}}
$$
Si l'application
$$
\alpha_j|_{T_i \times_T T_j} \circ
\sigma_{T_i \times_T T_j}^*\varphi_{ij} \circ
\alpha_i|_{T_i \times_T T_j} :
\mathcal{O}_{T_i \times_T T_j} \to \mathcal{O}_{T_i \times_T T_j}
$$
n'est pas la multiplication par $1$, mais par un élément $u_{ij}$, on peut multiplier
$\varphi_{ij}$ par $u_{ij}^{-1}$ pour y remédier. Ceci fait, considérons
l'endomorphisme
$$
\varphi_{ki}|_{X_{T_i \times_T T_j \times_T T_k}} \circ
\varphi_{jk}|_{X_{T_i \times_T T_j \times_T T_k}} \circ
\varphi_{ij}|_{X_{T_i \times_T T_j \times_T T_k}}
\quad\text{de}\quad
\mathcal{L}_i|_{X_{T_i \times_T T_j \times_T T_k}}
$$
qui est la multiplication par une fonction régulière $f_{ijk}$
sur le schéma $X_{T_i \times_T T_j \times_T T_k}$.
Par notre choix de $\varphi_{ij}$, l'image réciproque de
cette application par $\sigma$ est la multiplication par $1$. D'après
notre hypothèse sur les fonctions sur $X$, on voit que $f_{ijk} = 1$.
On obtient ainsi une donnée de descente pour le recouvrement fppf
$\{X_{T_i} \to X\}$. D'après
Descente, proposition \ref{descent-proposition-fpqc-descent-quasi-coherent},
il existe un $\mathcal{O}_{X_T}$-module inversible $\mathcal{L}$
et un isomorphisme $\alpha : \mathcal{O}_T \to \sigma_T^*\mathcal{L}$
dont l'image réciproque sur $X_{T_i}$ redonne $(\mathcal{L}_i, \alpha_i)$
(nous omettons un petit détail). Ainsi, $\mathcal{L}$ définit un objet
de $K(T)$, comme voulu.
\end{proof}



\section{Un critère de représentabilité}
\label{section-representability}

\noindent
Pour démontrer que le foncteur de Picard est représentable, nous utiliserons le
critère suivant.

\begin{lemma}
\label{lemma-criterion}
Soit $k$ un corps. Soit $G : (\Sch/k)^{opp} \to \textit{Groupes}$ un
foncteur. Avec la terminologie de
Schémas, définition \ref{schemes-definition-representable-by-open-immersions},
supposons que
\begin{enumerate}
\item $G$ soit un faisceau pour la topologie de Zariski,
\item il existe un sous-foncteur $F \subset G$ tel que
\begin{enumerate}
\item $F$ soit représentable,
\item $F \subset G$ soit représentable par une immersion ouverte,
\item pour toute extension de corps $K$ de $k$ et tout $g \in G(K)$,
il existe un $g' \in G(k)$ tel que $g'g \in F(K)$.
\end{enumerate}
\end{enumerate}
Alors $G$ est représentable par un schéma en groupes sur $k$.
\end{lemma}

\begin{proof}
Cela résulte de Schémas, lemme \ref{schemes-lemma-glue-functors}.
En effet, prenons $I = G(k)$ et, pour $i = g' \in I$, prenons pour $F_i \subset G$
le sous-foncteur qui associe à $T$ sur $k$ l'ensemble des éléments
$g \in G(T)$ tels que $g'g \in F(T)$. Alors $F_i \cong F$ par multiplication
par $g'$. L'application $F_i \to G$ est isomorphe à l'application $F \to G$
par multiplication par $g'$, donc est représentable par des immersions ouvertes.
Enfin, la famille $(F_i)_{i \in I}$ recouvre $G$ par l'hypothèse (2)(c).
Le lemme cité ci-dessus s'applique donc, ce qui achève la démonstration.
\end{proof}



\section{Le schéma de Picard d'une courbe}
\label{section-picard-curve}

\noindent
Dans cette section, nous appliquerons le lemme \ref{lemma-criterion} pour montrer que
$\Picardfunctor_{X/k}$ est représentable lorsque $k$ est un corps algébriquement
clos et $X$ une courbe projective lisse sur $k$. Pour cela,
nous utiliserons quelques résultats de cohomologie et de changement de base développés dans
le chapitre sur les catégories dérivées des schémas.

\begin{lemma}
\label{lemma-check-conditions}
Soit $k$ un corps. Soit $X$ une courbe projective lisse sur $k$
possédant un point $k$-rationnel. Alors les hypothèses du
lemme \ref{lemma-flat-geometrically-connected-fibres-with-section}
sont satisfaites.
\end{lemma}

\begin{proof}
L'expression « possède un point $k$-rationnel » signifie exactement que
le morphisme structural $f : X \to \Spec(k)$ admet une section, ce qui
vérifie la première condition.
D'après Variétés, lemme \ref{varieties-lemma-regular-functions-proper-variety},
le corps $k' = H^0(X, \mathcal{O}_X)$ est une extension de $k$.
Puisque $X$ possède un point $k$-rationnel, il existe un homomorphisme de $k$-algèbres
$k' \to k$, d'où $k' = k$.
Puisque $k$ est un corps, tout morphisme $T \to \Spec(k)$ est plat.
Par cohomologie et changement de base
(Cohomologie des schémas, lemme \ref{coherent-lemma-flat-base-change-cohomology}),
on voit donc que $\mathcal{O}_T \to f_{T, *}\mathcal{O}_{X_T}$ est un isomorphisme.
Cela achève la démonstration.
\end{proof}

\noindent
Soit $X$ une courbe projective lisse sur un corps $k$, munie d'un
point $k$-rationnel $\sigma$. Alors le foncteur
$$
\Picardfunctor_{X/k, \sigma} : (\Sch/k)^{opp} \longrightarrow \textit{Ab},\quad
T \longmapsto \Ker(\Pic(X_T) \xrightarrow{\sigma_T^*} \Pic(T))
$$
est isomorphe à $\Picardfunctor_{X/k}$ sur $(\Sch/k)_{fppf}$
d'après les lemmes \ref{lemma-check-conditions} et
\ref{lemma-flat-geometrically-connected-fibres-with-section}.
Il suffira donc de démontrer que $\Picardfunctor_{X/k, \sigma}$
est représentable. Nous emploierons la notation
« $\mathcal{L} \in \Picardfunctor_{X/k, \sigma}(T)$ » pour signifier que
$T$ est un schéma sur $k$ et que $\mathcal{L}$ est un
$\mathcal{O}_{X_T}$-module inversible dont la restriction à $T$ par $\sigma_T$
est isomorphe à $\mathcal{O}_T$.

\begin{lemma}
\label{lemma-define-open}
Soit $k$ un corps. Soit $X$ une courbe projective lisse sur $k$
munie d'un point $k$-rationnel $\sigma$. Pour tout schéma $T$ sur $k$,
considérons la partie $F(T) \subset \Picardfunctor_{X/k, \sigma}(T)$ formée des
$\mathcal{L}$ tels que $Rf_{T, *}\mathcal{L}$ soit isomorphe à un
$\mathcal{O}_T$-module inversible placé en degré $0$. Alors
$F \subset \Picardfunctor_{X/k, \sigma}$ est un sous-foncteur, et l'inclusion est
représentable par des immersions ouvertes.
\end{lemma}

\begin{proof}
Cela résulte immédiatement de Catégories dérivées des schémas, lemme
\ref{perfect-lemma-open-where-cohomology-in-degree-i-rank-r-geometric},
appliqué avec $i = 0$ et $r = 1$, et de
Schémas, définition \ref{schemes-definition-representable-by-open-immersions}.
\end{proof}

\noindent
Pour poursuivre, il convient de poser la définition suivante.

\begin{definition}
\label{definition-genus}
Soit $k$ un corps. Soit $X$ une courbe projective lisse géométriquement irréductible
sur $k$. Le {\it genre} de $X$ est $g = \dim_k H^1(X, \mathcal{O}_X)$.
\end{definition}

\begin{lemma}
\label{lemma-open-representable}
Soit $k$ un corps. Soit $X$ une courbe projective lisse de genre $g$
sur $k$, munie d'un point $k$-rationnel $\sigma$. Le sous-foncteur ouvert $F$ défini
au lemme \ref{lemma-define-open} est représentable par un sous-schéma ouvert de
$\underline{\Hilbfunctor}^g_{X/k}$.
\end{lemma}

\begin{proof}
Dans cette démonstration, les produits sans indice sont pris sur $\Spec(k)$.
D'après la proposition \ref{proposition-hilb-d}, le schéma
$H = \underline{\Hilbfunctor}^g_{X/k}$ existe.
Considérons le diviseur universel $D_{univ} \subset H \times X$
et le faisceau inversible associé $\mathcal{O}(D_{univ})$, voir
la remarque \ref{remark-universal-object-hilb-d}.
Nous le normalisons par produit tensoriel en utilisant l'image réciproque par
$\sigma_H : H \to H \times X$, pour obtenir
$$
\mathcal{L}_H =
\mathcal{O}(D_{univ})
\otimes_{\mathcal{O}_{H \times X}}
\text{pr}_H^*\sigma_H^*\mathcal{O}(D_{univ})^{\otimes -1}
\in
\Picardfunctor_{X/k, \sigma}(H)
$$
Par le lemme de Yoneda (Catégories, lemme \ref{categories-lemma-yoneda}),
le faisceau inversible $\mathcal{L}_H$ définit une transformation naturelle
$$
h_H \longrightarrow \Picardfunctor_{X/k, \sigma}
$$
Puisque $F$ est un sous-foncteur ouvert, il existe un plus grand ouvert
$W \subset H$ tel que $\mathcal{L}_H|_{W \times X}$ appartienne à
$F(W)$. Cet ouvert n'est autre que le
sous-schéma ouvert construit dans
Catégories dérivées des schémas, lemme
\ref{perfect-lemma-open-where-cohomology-in-degree-i-rank-r-geometric},
avec $i = 0$ et $r = 1$, pour le morphisme $H \times X \to H$ et le faisceau
$\mathcal{F} = \mathcal{O}(D_{univ})$. En appliquant à nouveau le lemme de Yoneda,
on obtient un diagramme commutatif
$$
\xymatrix{
h_W \ar[d] \ar[r] & F \ar[d] \\
h_H \ar[r] & \Picardfunctor_{X/k, \sigma}
}
$$
Pour achever la démonstration, nous montrerons que la flèche horizontale supérieure est un
isomorphisme.

\medskip\noindent
Soit $\mathcal{L} \in F(T) \subset \Picardfunctor_{X/k, \sigma}(T)$.
Soit $\mathcal{N}$ le $\mathcal{O}_T$-module inversible
tel que $Rf_{T, *}\mathcal{L} \cong \mathcal{N}[0]$.
L'application d'adjonction
$$
f_T^*\mathcal{N} \longrightarrow \mathcal{L}
\quad\text{correspond à une section }s\text{ de}\quad
\mathcal{L} \otimes f_T^*\mathcal{N}^{\otimes -1}
$$
sur $X_T$. Assertion : le schéma des zéros de $s$ est un diviseur de Cartier effectif
relatif $D$ sur $(T \times X)/T$, fini localement libre de degré $g$ sur $T$.

\medskip\noindent
Achevons la démonstration du lemme en admettant cette assertion.
Le diviseur $D$ définit un morphisme $m : T \to H$ tel que $D$ soit l'image réciproque de
$D_{univ}$. Alors
$$
(m \times \text{id}_X)^*\mathcal{O}(D_{univ}) \cong
\mathcal{O}_{T \times X}(D)
$$
Puisque $\mathcal{O}_{T \times X}(D) \cong
\mathcal{L} \otimes f_T^*\mathcal{N}^{\otimes -1}$,
on en conclut que $(m \times \text{id}_X)^*\mathcal{L}_H$ et $\mathcal{L}$
diffèrent par tensorisation avec l'image réciproque d'un module inversible sur $T$.
Cela montre que $m : T \to H$ se factorise par l'ouvert $W \subset H$ ci-dessus.
De plus, il s'ensuit que ces modules inversibles définissent, après normalisation
par image réciproque suivant $\sigma_T$ comme ci-dessus, le même élément de
$\Picardfunctor_{X/k, \sigma}(T)$. Une chasse au diagramme utilisant le lemme de Yoneda
montre que $m \in h_W(T)$ a pour image $\mathcal{L} \in F(T)$. Nous omettons
la vérification du fait que la règle $F(T) \to h_W(T)$,
$\mathcal{L} \mapsto m$, définit un inverse de la transformation
de foncteurs ci-dessus.

\medskip\noindent
Démonstration de l'assertion. Puisque $D$ est un sous-schéma fermé localement principal
de $T \times X$, il suffit de montrer que les fibres de $D$ sur $T$ sont des
diviseurs de Cartier effectifs, voir le lemme \ref{lemma-divisors-on-curves} et
Diviseurs, lemme \ref{divisors-lemma-fibre-Cartier}. Puisque le calcul
de la cohomologie de $\mathcal{L}$ commute au changement de base
(Catégories dérivées des schémas, lemme
\ref{perfect-lemma-flat-proper-perfect-direct-image-general}),
on se ramène à $T = \Spec(K)$, où $K/k$ est une extension de corps.
Alors $\mathcal{L}$ est un faisceau inversible sur $X_K$ tel que
$H^0(X_K, \mathcal{L}) = K$ et $H^1(X_K, \mathcal{L}) = 0$. Ainsi,
$$
\deg(\mathcal{L}) = \chi(X_K, \mathcal{L}) - \chi(X_K, \mathcal{O}_{X_K})
= 1 - (1 - g) = g
$$
Voir Variétés, définition \ref{varieties-definition-degree-invertible-sheaf}.
Pour achever la démonstration, il faut montrer qu'une section non nulle de $\mathcal{L}$
définit un diviseur de Cartier effectif sur $X_K$.
C'est clair.
\end{proof}

\begin{lemma}
\label{lemma-twist-with-general-divisor}
Soit $k$ un corps séparablement clos. Soit $X$ une courbe projective lisse
de genre $g$ sur $k$. Soit $K/k$ une extension de corps et soit
$\mathcal{L}$ un faisceau inversible sur $X_K$. Il existe alors un
faisceau inversible $\mathcal{L}_0$ sur $X$ tel que
$\dim_K H^0(X_K,
\mathcal{L} \otimes_{\mathcal{O}_{X_K}} \mathcal{L}_0|_{X_K}) = 1$ et
$\dim_K H^1(X_K,
\mathcal{L} \otimes_{\mathcal{O}_{X_K}} \mathcal{L}_0|_{X_K}) = 0$.
\end{lemma}

\begin{proof}
Cette démonstration est une variante de celle de
Variétés, lemme \ref{varieties-lemma-general-degree-g-line-bundle}.
Nous invitons le lecteur à lire d'abord cette dernière.

\medskip\noindent
Choisissons d'abord un faisceau inversible ample $\mathcal{L}_0$ et
remplaçons $\mathcal{L}$ par
$\mathcal{L} \otimes_{\mathcal{O}_{X_K}} \mathcal{L}_0^{\otimes n}|_{X_K}$
pour un $n \gg 0$. Nous pourrons ainsi supposer que
$H^0(X_K, \mathcal{L}) \not = 0$ et $H^1(X_K, \mathcal{L}) = 0$.
En effet, l'annulation résulte de Cohomologie des schémas, lemme
\ref{coherent-lemma-vanshing-gives-ample}, et la non-annulation de ce que
le degré du produit tensoriel est $\gg 0$.
Nous achèverons la démonstration par récurrence descendante sur
$t = \dim_K H^0(X_K, \mathcal{L})$. Le cas initial $t = 1$ est trivial.
Supposons $t > 1$.

\medskip\noindent
Notons que, pour un point $k$-rationnel $x$ de $X$, l'image réciproque $x_K$
est un point $K$-rationnel de $X_K$. De plus, il existe une infinité de
points $k$-rationnels d'après Variétés, lemme
\ref{varieties-lemma-smooth-separable-closed-points-dense}. Par conséquent,
les points $x_K$ forment une famille Zariski-dense de points de $X_K$.

\medskip\noindent
Soit $s \in H^0(X_K, \mathcal{L})$ non nul. Le paragraphe précédent
montre qu'il existe un point $k$-rationnel
$x$ tel que $s$ ne s'annule pas en $x_K$. Soit $\mathcal{I}$
le faisceau d'idéaux de $i : x_K \to X_K$ comme dans
Variétés, lemme \ref{varieties-lemma-regular-point-on-curve}. Considérons la
suite exacte courte
$$
0 \to \mathcal{I} \otimes_{\mathcal{O}_{X_K}} \mathcal{L} \to
\mathcal{L} \to i_*i^*\mathcal{L} \to 0
$$
Notons que $H^0(X_K, i_*i^*\mathcal{L}) = H^0(x_K, i^*\mathcal{L})$
est de dimension $1$ sur $K$. Puisque $s$ ne s'annule pas en $x$, on en conclut que
$$
H^0(X_K, \mathcal{L}) \longrightarrow H^0(X, i_*i^*\mathcal{L})
$$
est surjective. Ainsi,
$\dim_K H^0(X_K, \mathcal{I} \otimes_{\mathcal{O}_{X_K}} \mathcal{L}) = t - 1$.
Enfin, la suite exacte longue de cohomologie montre aussi que
$H^1(X_K, \mathcal{I} \otimes_{\mathcal{O}_{X_K}} \mathcal{L}) = 0$,
ce qui achève l'étape de récurrence.
\end{proof}

\begin{proposition}
\label{proposition-pic-curve}
Soit $k$ un corps séparablement clos. Soit $X$ une courbe projective lisse
sur $k$. Le foncteur de Picard $\Picardfunctor_{X/k}$ est représentable.
\end{proposition}

\begin{proof}
Puisque $k$ est séparablement clos, il existe un point $k$-rationnel $\sigma$
de $X$, voir
Variétés, lemme \ref{varieties-lemma-smooth-separable-closed-points-dense}.
Comme on l'a vu ci-dessus, il suffit de montrer que le foncteur
$\Picardfunctor_{X/k, \sigma}$ classifiant les modules inversibles triviaux le long de
$\sigma$ est représentable. Pour cela, nous vérifierons les conditions (1),
(2)(a), (2)(b) et (2)(c) du
lemme \ref{lemma-criterion}.

\medskip\noindent
Le foncteur $\Picardfunctor_{X/k, \sigma}$ satisfait à la condition de faisceau
pour la topologie fppf, car il est isomorphe à $\Picardfunctor_{X/k}$.
Il serait plus exact de dire que nous avons démontré la condition de faisceau
pour $\Picardfunctor_{X/k, \sigma}$ dans la démonstration du
lemme \ref{lemma-flat-geometrically-connected-fibres-with-section},
qui s'applique d'après le lemme \ref{lemma-check-conditions}.
Cela démontre la condition (1).

\medskip\noindent
Prenons pour sous-foncteur $F$ celui du lemme \ref{lemma-define-open}.
La condition (2)(b) en résulte.
La condition (2)(a) est le lemme \ref{lemma-open-representable}.
La condition (2)(c) est le lemme \ref{lemma-twist-with-general-divisor}.
\end{proof}

\noindent
La démonstration précédente donne en fait des renseignements supplémentaires, que nous
rassemblons ici.

\begin{lemma}
\label{lemma-picard-pieces}
Soit $k$ un corps séparablement clos. Soit $X$ une courbe projective lisse
de genre $g$ sur $k$.
\begin{enumerate}
\item $\underline{\Picardfunctor}_{X/k}$ est une réunion disjointe de variétés
propres et lisses de dimension $g$, notées $\underline{\Picardfunctor}^d_{X/k}$,
\item les points à valeurs dans $k$ de $\underline{\Picardfunctor}^d_{X/k}$
correspondent aux $\mathcal{O}_X$-modules inversibles de degré $d$,
\item $\underline{\Picardfunctor}^0_{X/k}$
est un sous-schéma en groupes ouvert et fermé,
\item pour $d \geq 0$, il existe un morphisme canonique
$\gamma_d :
\underline{\Hilbfunctor}^d_{X/k} \to \underline{\Picardfunctor}^d_{X/k}$
\item les morphismes $\gamma_d$
sont surjectifs pour $d \geq g$ et lisses pour $d \geq 2g - 1$,
\item le morphisme
$\underline{\Hilbfunctor}^g_{X/k} \to \underline{\Picardfunctor}^g_{X/k}$
est birationnel.
\end{enumerate}
\end{lemma}

\begin{proof}
Choisissons un point $k$-rationnel $\sigma$ de $X$. Rappelons que $\Picardfunctor_{X/k}$
est isomorphe au foncteur $\Picardfunctor_{X/k, \sigma}$. D'après
Catégories dérivées des schémas, lemme
\ref{perfect-lemma-chi-locally-constant-geometric},
pour tout $d \in \mathbf{Z}$, il existe un sous-foncteur ouvert
$$
\Picardfunctor^d_{X/k, \sigma} \subset \Picardfunctor_{X/k, \sigma}
$$
dont la valeur sur un schéma $T$ sur $k$ est formée des
$\mathcal{L} \in \Picardfunctor_{X/k, \sigma}(T)$ tels que
$\chi(X_t, \mathcal{L}_t) = d + 1 - g$, et l'on a en outre
$$
\Picardfunctor_{X/k, \sigma} =
\coprod\nolimits_{d \in \mathbf{Z}} \Picardfunctor^d_{X/k, \sigma}
$$
en tant que faisceaux fppf. Le schéma $\underline{\Picardfunctor}_{X/k}$
(qui existe d'après la proposition \ref{proposition-pic-curve})
admet donc une décomposition correspondante
$$
\underline{\Picardfunctor}_{X/k, \sigma} =
\coprod\nolimits_{d \in \mathbf{Z}} \underline{\Picardfunctor}^d_{X/k, \sigma}
$$
où les points de $\underline{\Picardfunctor}^d_{X/k, \sigma}$ correspondent
aux classes d'isomorphisme de modules inversibles de degré $d$ sur $X$.

\medskip\noindent
Fixons $d \geq 0$. Il existe un morphisme
$$
\gamma_d :
\underline{\Hilbfunctor}^d_{X/k}
\longrightarrow
\underline{\Picardfunctor}^d_{X/k}
$$
provenant du faisceau inversible $\mathcal{O}(D_{univ})$ sur
$\underline{\Hilbfunctor}^d_{X/k} \times_k X$
(remarque \ref{remark-universal-object-hilb-d}) par le lemme de Yoneda
(Catégories, lemme \ref{categories-lemma-yoneda}).
Notre démonstration de la représentabilité du foncteur de Picard de $X/k$
dans la proposition \ref{proposition-pic-curve} et le
lemme \ref{lemma-open-representable} montre que $\gamma_g$
induit une immersion ouverte sur un ouvert non vide de
$\underline{\Hilbfunctor}^g_{X/k}$. De plus, la démonstration montre
que les translatés de cet ouvert par les points $k$-rationnels du
schéma en groupes $\underline{\Picardfunctor}_{X/k}$ définissent
un recouvrement ouvert. Puisque
$\underline{\Hilbfunctor}^g_{X/K}$ est lisse de dimension $g$
(proposition \ref{proposition-hilb-d})
sur $k$, on en conclut que le
schéma en groupes $\underline{\Picardfunctor}_{X/k}$ est lisse de dimension $g$
sur $k$.

\medskip\noindent
D'après Groupoïdes, lemme \ref{groupoids-lemma-group-scheme-over-field-separated},
le schéma $\underline{\Picardfunctor}_{X/k}$ est séparé.
Ainsi, pour tout $d \geq 0$, l'image de $\gamma_d$
est une variété propre sur $k$
(Morphismes, lemme \ref{morphisms-lemma-scheme-theoretic-image-is-proper}).

\medskip\noindent
Soit $d \geq g$. Pour toute extension de corps $K/k$ et tout
$\mathcal{O}_{X_K}$-module inversible $\mathcal{L}$ de degré $d$, on a
$\chi(X_K, \mathcal{L}) = d + 1 - g > 0$. Ainsi, $\mathcal{L}$ admet une
section non nulle, et l'on en conclut que $\mathcal{L} = \mathcal{O}_{X_K}(D)$
pour un diviseur $D \subset X_K$ de degré $d$. Il s'ensuit que
$\gamma_d$ est surjectif.

\medskip\noindent
En combinant les faits précédents, on voit que
$\underline{\Picardfunctor}^d_{X/k}$ est propre pour $d \geq g$.
Cela achève la démonstration de (2), car nous savons maintenant que
$\underline{\Picardfunctor}^d_{X/k}$ est propre pour $d \geq g$, et
tous les $\underline{\Picardfunctor}^d_{X/k}$ sont alors propres par translation.

\medskip\noindent
Il reste à montrer que $\gamma_d$ est lisse pour $d \geq 2g - 1$.
Considérons un $\mathcal{O}_X$-module inversible $\mathcal{L}$ de degré
$d$. La fibre au-dessus du point correspondant à $\mathcal{L}$ est alors
$$
Z = \{D \subset X \mid \mathcal{O}_X(D) \cong \mathcal{L}\} \subset
\underline{\Hilbfunctor}^d_{X/k}
$$
munie de sa structure naturelle de schéma. Puisque tout isomorphisme
$\mathcal{O}_X(D) \to \mathcal{L}$ est bien défini à
multiplication près par un scalaire non nul, la section canonique
$1 \in \mathcal{O}_X(D)$ a pour image une section
$s \in \Gamma(X, \mathcal{L})$ bien définie à multiplication
près par un scalaire non nul. On obtient ainsi un morphisme
$$
Z \longrightarrow
\text{Proj}(\text{Sym}(\Gamma(X, \mathcal{L})^*))
$$
(le dual intervient en raison de nos conventions). Ce morphisme est un isomorphisme,
car, étant donné une section de $\mathcal{L}$, on peut prendre le diviseur de Cartier
effectif associé ; autrement dit, on peut construire un inverse du
morphisme écrit ci-dessus. Nous omettons la formulation précise et la démonstration.
Puisque $\dim H^0(X, \mathcal{L}) = d + 1 - g$ pour tout
$\mathcal{L}$ de degré $d \geq 2g - 1$ d'après
Variétés, lemme \ref{varieties-lemma-vanishing-degree-2g-and-1-line-bundle},
on voit que $\text{Proj}(\text{Sym}(\Gamma(X, \mathcal{L})^*))
\cong \mathbf{P}^{d - g}_k$.
On en conclut que $\dim(Z) = \dim(\mathbf{P}^{d - g}_k) = d - g$.
Ainsi, les fibres du morphisme $\gamma_d$ sont toutes
de dimension égale à la différence des dimensions de
$\underline{\Hilbfunctor}^d_{X/k}$ et de $\underline{\Picardfunctor}^d_{X/k}$.
Il s'ensuit que $\gamma_d$ est plat, voir
Algèbre, lemme \ref{algebra-lemma-CM-over-regular-flat}.
Comme ses fibres sont en outre lisses, on en conclut que $\gamma_d$
est lisse d'après Morphismes, lemme \ref{morphisms-lemma-smooth-flat-smooth-fibres}.
\end{proof}





\section{Quelques remarques sur les groupes de Picard}
\label{section-remarks-picard}

\noindent
Cette section poursuit la discussion de
Variétés, section \ref{varieties-section-picard-groups},
qui sera prolongée dans
Courbes algébriques, section \ref{curves-section-torsion-in-pic}.

\begin{lemma}
\label{lemma-pic-descends}
Soit $k$ un corps. Soit $X$ un schéma quasi-compact et quasi-séparé
sur $k$ tel que $H^0(X, \mathcal{O}_X) = k$. Si $X$ possède un point $k$-rationnel,
alors, pour toute extension galoisienne $k'/k$, on a
$$
\Pic(X) = \Pic(X_{k'})^{\text{Gal}(k'/k)}
$$
De plus, l'action de $\text{Gal}(k'/k)$ sur $\Pic(X_{k'})$
est continue.
\end{lemma}

\begin{proof}
Puisque $\text{Gal}(k'/k) = \text{Aut}(k'/k)$, ce groupe opère (à droite)
sur $\Spec(k')$, donc opère (à droite) sur
$X_{k'} = X \times_{\Spec(k)} \Spec(k')$ ; comme $\Pic(-)$
est un foncteur contravariant, il opère (à gauche) sur $\Pic(X_{k'})$.
Si $k'/k$ est une extension galoisienne infinie, on écrit
$k' = \colim k'_\lambda$ comme limite inductive filtrante d'extensions galoisiennes
finies, voir Corps, lemme \ref{fields-lemma-infinite-galois-limit}.
Alors $X_{k'} = \lim X_{k_\lambda}$
(comme dans Limites, section \ref{limits-section-limits})
et l'on obtient
$$
\Pic(X_{k'}) = \colim \Pic(X_{k_\lambda})
$$
d'après Limites, lemme \ref{limits-lemma-descend-invertible-modules}.
De plus, les applications de transition de ce système de groupes abéliens
sont injectives d'après Variétés, lemme \ref{varieties-lemma-change-fields-pic}.
Il s'ensuit que tout élément de $\Pic(X_{k'})$ est fixé par
l'un des sous-groupes ouverts $\text{Gal}(k'/k_\lambda)$, ce qui signifie exactement
que l'action est continue. L'injectivité des applications de transition
montre qu'il suffit de démontrer l'assertion sur les points fixes dans le
cas où $k'/k$ est une extension galoisienne finie.

\medskip\noindent
Supposons $k'/k$ galoisienne finie, de groupe de Galois $G = \text{Gal}(k'/k)$.
Soit $\mathcal{L}$ un élément de $\Pic(X_{k'})$ fixé par $G$.
Nous utiliserons la descente galoisienne (Descente, lemme \ref{descent-lemma-galois-descent})
pour montrer que $\mathcal{L}$ est l'image réciproque d'un faisceau inversible sur $X$.
Rappelons que $f_\sigma = \text{id}_X \times \Spec(\sigma) : X_{k'} \to X_{k'}$
et que $\sigma$ opère sur $\Pic(X_{k'})$ par image réciproque suivant $f_\sigma$.
Pour tout $\sigma \in G$, on peut donc choisir un isomorphisme
$\varphi_\sigma : \mathcal{L} \to f_\sigma^*\mathcal{L}$,
puisque $\mathcal{L}$ est fixé par l'action de $G$.
La difficulté est que l'on ne sait pas si l'on peut choisir
$\varphi_\sigma$ de manière à satisfaire à la condition de cocycle
$\varphi_{\sigma\tau} = f_\sigma^*\varphi_\tau \circ \varphi_\sigma$.
Pour voir que c'est possible, nous utilisons le fait que $X$ possède un point $k$-rationnel
$x \in X(k)$. Bien entendu, $x$ détermine de même un point $k'$-rationnel
$x' \in X_{k'}$ fixé par $f_\sigma$ pour tout $\sigma$.
Choisissons un élément non nul $s$ de la fibre de $\mathcal{L}$ en $x'$ ;
cette fibre est l'espace vectoriel de dimension $1$ sur $k' = \kappa(x')$
$$
\mathcal{L}_{x'} \otimes_{\mathcal{O}_{X_{k'}, x'}} \kappa(x').
$$
Alors $f_\sigma^*s$ est un élément non nul de la fibre de
$f_\sigma^*\mathcal{L}$ en $x'$. Puisqu'on peut multiplier $\varphi_\sigma$
par un élément de $(k')^*$, on peut supposer que $\varphi_\sigma$ envoie
$s$ sur $f_\sigma^*s$. On voit alors que $\varphi_{\sigma\tau}$ et
$f_\sigma^*\varphi_\tau \circ \varphi_\sigma$
envoient tous deux $s$ sur $f_{\sigma\tau}^*s = f_\tau^*f_\sigma^*s$.
Puisque $H^0(X_{k'}, \mathcal{O}_{X_{k'}}) = k'$, ces deux isomorphismes
sont nécessairement égaux (l'un est le produit de l'autre par une unité globale
et ils coïncident en $x'$), ce qui achève la démonstration.
\end{proof}

\begin{lemma}
\label{lemma-torsion-descends}
Soit $k$ un corps de caractéristique $p > 0$. Soit $X$ un
schéma quasi-compact et quasi-séparé sur $k$ tel que
$H^0(X, \mathcal{O}_X) = k$. Soit $n$ un entier premier à $p$.
Alors l'application
$$
\Pic(X)[n] \longrightarrow \Pic(X_{k'})[n]
$$
est bijective pour toute extension purement inséparable $k'/k$.
\end{lemma}

\begin{proof}
Notons d'abord que l'application $\Pic(X) \to \Pic(X_{k'})$
est injective d'après Variétés, lemme \ref{varieties-lemma-change-fields-pic}.
Il faut donc montrer que l'application du lemme est surjective.
Soit $\mathcal{L}$ un $\mathcal{O}_{X_{k'}}$-module inversible
dont la classe est d'ordre divisant $n$ dans $\Pic(X_{k'})$.
Choisissons un isomorphisme
$\alpha : \mathcal{L}^{\otimes n} \to \mathcal{O}_{X_{k'}}$
de modules inversibles.
Nous allons montrer que le couple $(\mathcal{L}, \alpha)$ descend
à $X$.

\medskip\noindent
Posons $A = k' \otimes_k k'$. Puisque $k'/k$ est purement inséparable,
le noyau de l'application de multiplication $A \to k'$ est un
idéal localement nilpotent $I$ de $A$. Notons que
$$
X_A = X \times_{\Spec(k)} \Spec(A) = X_{k'} \times_X X_{k'}
$$
est muni de deux projections $\text{pr}_i : X_A \to X_{k'}$, $i = 0, 1$,
qui coïncident sur $A/I$. Les modules inversibles
$\mathcal{L}_i = \text{pr}_i^*\mathcal{L}$
coïncident donc sur le sous-schéma fermé $X_{A/I} = X_{k'}$.
Puisque $X_{A/I} \to X_A$ est un épaississement et que
les $\mathcal{L}_i$ sont de $n$-torsion, il
existe un isomorphisme $\varphi : \mathcal{L}_0 \to \mathcal{L}_1$ d'après
Compléments sur les morphismes, lemme \ref{more-morphisms-lemma-torsion-pic-thickening}.
On peut choisir $\varphi$ dont la réduction modulo $I$ soit l'identité. En effet,
$H^0(X, \mathcal{O}_X) = k$ entraîne
$H^0(X_{k'}, \mathcal{O}_{X_{k'}}) = k'$ d'après
Cohomologie des schémas, lemme \ref{coherent-lemma-flat-base-change-cohomology},
et $A \to k'$ est surjective ; on peut donc ajuster $\varphi$ en le multipliant par
un élément convenable de $A$. Considérons l'application
$$
\lambda : \mathcal{O}_{X_A}
\xrightarrow{\text{pr}_0^*\alpha^{-1}}
\mathcal{L}_0^{\otimes n}
\xrightarrow{\varphi^{\otimes n}}
\mathcal{L}_1^{\otimes n}
\xrightarrow{\text{pr}_0^*\alpha}
\mathcal{O}_{X_A}
$$
On peut regarder $\lambda$ comme un élément de $A$, car
$H^0(X_A, \mathcal{O}_{X_A}) = A$ (même référence que ci-dessus).
Puisque $\varphi$ se réduit à l'identité modulo $I$, on
voit que $\lambda = 1 \bmod I$. Il existe alors une unique
racine $n$-ième de $\lambda$ dans $1 + I$
(Algèbre, lemme \ref{algebra-lemma-lift-nth-roots})
et, en multipliant $\varphi$ par son inverse, on obtient $\lambda = 1$.
Nous affirmons que $(\mathcal{L}, \varphi)$ est une donnée de descente
pour le recouvrement fpqc $\{X_{k'} \to X\}$
(Descente, définition \ref{descent-definition-descent-datum-quasi-coherent}).
Si cette assertion est vraie, $\mathcal{L}$ est l'image réciproque d'un
$\mathcal{O}_X$-module inversible $\mathcal{N}$ d'après
Descente, proposition \ref{descent-proposition-fpqc-descent-quasi-coherent}.
L'injectivité de l'application entre les groupes de Picard montre que $\mathcal{N}$
est un élément de torsion de $\Pic(X)$ de même ordre que $\mathcal{L}$.

\medskip\noindent
Démonstration de l'assertion. Il faut vérifier que
$$
\text{pr}_{12}^*\varphi \circ
\text{pr}_{01}^*\varphi =
\text{pr}_{02}^*\varphi
\quad\text{sur}\quad
X_{k'} \times_X X_{k'} \times_X X_{k'} = X_{k' \otimes_k k' \otimes_k k'}
$$
Comme précédemment, le morphisme diagonal
$\Delta : X_{k'} \to X_{k' \otimes_k k' \otimes_k k'}$
est un épaississement. Les membres de gauche et de droite de l'égalité
sont des applications $a, b : p_0^*\mathcal{L} \to p_2^*\mathcal{L}$ compatibles avec
$p_0^*\alpha$ et $p_2^*\alpha$, où
$p_i : X_{k' \otimes_k k' \otimes_k k'} \to X_{k'}$
sont les morphismes de projection. Enfin, les images réciproques de $a, b$ par
$\Delta$ sont égales.
Localement sur des ouverts affines (dans des trivialisations locales), cela signifie que
$a, b$ sont données par multiplication par des fonctions inversibles
qui se réduisent à la même fonction modulo un idéal localement nilpotent
et dont les puissances $n$-ièmes sont égales. Il résulte alors de
Algèbre, lemme \ref{algebra-lemma-lift-nth-roots}
que ces fonctions sont égales.
\end{proof}





\begin{multicols}{2}[\section{Autres chapitres}]
\noindent
Préliminaires
\begin{enumerate}
\item \hyperref[introduction-section-phantom]{Introduction}
\item \hyperref[conventions-section-phantom]{Conventions}
\item \hyperref[sets-section-phantom]{Théorie des ensembles}
\item \hyperref[categories-section-phantom]{Catégories}
\item \hyperref[topology-section-phantom]{Topologie}
\item \hyperref[sheaves-section-phantom]{Faisceaux sur les espaces}
\item \hyperref[sites-section-phantom]{Sites et faisceaux}
\item \hyperref[stacks-section-phantom]{Champs}
\item \hyperref[fields-section-phantom]{Corps}
\item \hyperref[algebra-section-phantom]{Algèbre commutative}
\item \hyperref[brauer-section-phantom]{Groupes de Brauer}
\item \hyperref[homology-section-phantom]{Algèbre homologique}
\item \hyperref[derived-section-phantom]{Catégories dérivées}
\item \hyperref[simplicial-section-phantom]{Méthodes simpliciales}
\item \hyperref[more-algebra-section-phantom]{Compléments d'algèbre}
\item \hyperref[smoothing-section-phantom]{Lissage des morphismes d'anneaux}
\item \hyperref[modules-section-phantom]{Faisceaux de modules}
\item \hyperref[sites-modules-section-phantom]{Modules sur les sites}
\item \hyperref[injectives-section-phantom]{Objets injectifs}
\item \hyperref[cohomology-section-phantom]{Cohomologie des faisceaux}
\item \hyperref[sites-cohomology-section-phantom]{Cohomologie sur les sites}
\item \hyperref[dga-section-phantom]{Algèbre différentielle graduée}
\item \hyperref[dpa-section-phantom]{Algèbre à puissances divisées}
\item \hyperref[sdga-section-phantom]{Faisceaux différentiels gradués}
\item \hyperref[hypercovering-section-phantom]{Hyperrecouvrements}
\end{enumerate}
Schémas
\begin{enumerate}
\setcounter{enumi}{25}
\item \hyperref[schemes-section-phantom]{Schémas}
\item \hyperref[constructions-section-phantom]{Constructions de schémas}
\item \hyperref[properties-section-phantom]{Propriétés des schémas}
\item \hyperref[morphisms-section-phantom]{Morphismes de schémas}
\item \hyperref[coherent-section-phantom]{Cohomologie des schémas}
\item \hyperref[divisors-section-phantom]{Diviseurs}
\item \hyperref[limits-section-phantom]{Limites de schémas}
\item \hyperref[varieties-section-phantom]{Variétés}
\item \hyperref[topologies-section-phantom]{Topologies sur les schémas}
\item \hyperref[descent-section-phantom]{Descente}
\item \hyperref[perfect-section-phantom]{Catégories dérivées de schémas}
\item \hyperref[more-morphisms-section-phantom]{Compléments sur les morphismes}
\item \hyperref[flat-section-phantom]{Compléments sur la platitude}
\item \hyperref[groupoids-section-phantom]{Schémas en groupoïdes}
\item \hyperref[more-groupoids-section-phantom]{Compléments sur les schémas en groupoïdes}
\item \hyperref[etale-section-phantom]{Morphismes étales de schémas}
\end{enumerate}
Thèmes de théorie des schémas
\begin{enumerate}
\setcounter{enumi}{41}
\item \hyperref[chow-section-phantom]{Homologie de Chow}
\item \hyperref[intersection-section-phantom]{Théorie de l'intersection}
\item \hyperref[pic-section-phantom]{Schémas de Picard des courbes}
\item \hyperref[weil-section-phantom]{Théories cohomologiques de Weil}
\item \hyperref[adequate-section-phantom]{Modules adéquats}
\item \hyperref[dualizing-section-phantom]{Complexes dualisants}
\item \hyperref[duality-section-phantom]{Dualité pour les schémas}
\item \hyperref[discriminant-section-phantom]{Discriminants et différentes}
\item \hyperref[derham-section-phantom]{Cohomologie de de Rham}
\item \hyperref[local-cohomology-section-phantom]{Cohomologie locale}
\item \hyperref[algebraization-section-phantom]{Géométrie algébrique et formelle}
\item \hyperref[curves-section-phantom]{Courbes algébriques}
\item \hyperref[resolve-section-phantom]{Résolution des surfaces}
\item \hyperref[models-section-phantom]{Réduction semi-stable}
\item \hyperref[functors-section-phantom]{Foncteurs et morphismes}
\item \hyperref[equiv-section-phantom]{Catégories dérivées de variétés}
\item \hyperref[pione-section-phantom]{Groupes fondamentaux de schémas}
\item \hyperref[etale-cohomology-section-phantom]{Cohomologie étale}
\item \hyperref[crystalline-section-phantom]{Cohomologie cristalline}
\item \hyperref[proetale-section-phantom]{Cohomologie pro-étale}
\item \hyperref[relative-cycles-section-phantom]{Cycles relatifs}
\item \hyperref[more-etale-section-phantom]{Compléments de cohomologie étale}
\item \hyperref[trace-section-phantom]{La formule des traces}
\end{enumerate}
Espaces algébriques
\begin{enumerate}
\setcounter{enumi}{64}
\item \hyperref[spaces-section-phantom]{Espaces algébriques}
\item \hyperref[spaces-properties-section-phantom]{Propriétés des espaces algébriques}
\item \hyperref[spaces-morphisms-section-phantom]{Morphismes d'espaces algébriques}
\item \hyperref[decent-spaces-section-phantom]{Espaces algébriques décents}
\item \hyperref[spaces-cohomology-section-phantom]{Cohomologie des espaces algébriques}
\item \hyperref[spaces-limits-section-phantom]{Limites d'espaces algébriques}
\item \hyperref[spaces-divisors-section-phantom]{Diviseurs sur les espaces algébriques}
\item \hyperref[spaces-over-fields-section-phantom]{Espaces algébriques sur des corps}
\item \hyperref[spaces-topologies-section-phantom]{Topologies sur les espaces algébriques}
\item \hyperref[spaces-descent-section-phantom]{Descente et espaces algébriques}
\item \hyperref[spaces-perfect-section-phantom]{Catégories dérivées d'espaces}
\item \hyperref[spaces-more-morphisms-section-phantom]{Compléments sur les morphismes d'espaces}
\item \hyperref[spaces-flat-section-phantom]{Platitude sur les espaces algébriques}
\item \hyperref[spaces-groupoids-section-phantom]{Groupoïdes dans les espaces algébriques}
\item \hyperref[spaces-more-groupoids-section-phantom]{Compléments sur les groupoïdes dans les espaces}
\item \hyperref[bootstrap-section-phantom]{Amorçage}
\item \hyperref[spaces-pushouts-section-phantom]{Sommes amalgamées d'espaces algébriques}
\end{enumerate}
Thèmes de géométrie
\begin{enumerate}
\setcounter{enumi}{81}
\item \hyperref[spaces-chow-section-phantom]{Groupes de Chow des espaces}
\item \hyperref[groupoids-quotients-section-phantom]{Quotients de groupoïdes}
\item \hyperref[spaces-more-cohomology-section-phantom]{Compléments sur la cohomologie des espaces}
\item \hyperref[spaces-simplicial-section-phantom]{Espaces simpliciaux}
\item \hyperref[spaces-duality-section-phantom]{Dualité pour les espaces}
\item \hyperref[formal-spaces-section-phantom]{Espaces algébriques formels}
\item \hyperref[restricted-section-phantom]{Algébrisation des espaces formels}
\item \hyperref[spaces-resolve-section-phantom]{Résolution des surfaces revisitée}
\end{enumerate}
Théorie des déformations
\begin{enumerate}
\setcounter{enumi}{89}
\item \hyperref[formal-defos-section-phantom]{Théorie formelle des déformations}
\item \hyperref[defos-section-phantom]{Théorie des déformations}
\item \hyperref[cotangent-section-phantom]{Le complexe cotangent}
\item \hyperref[examples-defos-section-phantom]{Problèmes de déformation}
\end{enumerate}
Champs algébriques
\begin{enumerate}
\setcounter{enumi}{93}
\item \hyperref[algebraic-section-phantom]{Champs algébriques}
\item \hyperref[examples-stacks-section-phantom]{Exemples de champs}
\item \hyperref[stacks-sheaves-section-phantom]{Faisceaux sur les champs algébriques}
\item \hyperref[criteria-section-phantom]{Critères de représentabilité}
\item \hyperref[artin-section-phantom]{Axiomes d'Artin}
\item \hyperref[quot-section-phantom]{Espaces de Quot et de Hilbert}
\item \hyperref[stacks-properties-section-phantom]{Propriétés des champs algébriques}
\item \hyperref[stacks-morphisms-section-phantom]{Morphismes de champs algébriques}
\item \hyperref[stacks-limits-section-phantom]{Limites de champs algébriques}
\item \hyperref[stacks-cohomology-section-phantom]{Cohomologie des champs algébriques}
\item \hyperref[stacks-perfect-section-phantom]{Catégories dérivées de champs}
\item \hyperref[stacks-introduction-section-phantom]{Introduction aux champs algébriques}
\item \hyperref[stacks-more-morphisms-section-phantom]{Compléments sur les morphismes de champs}
\item \hyperref[stacks-geometry-section-phantom]{La géométrie des champs}
\end{enumerate}
Thèmes de théorie des espaces de modules
\begin{enumerate}
\setcounter{enumi}{107}
\item \hyperref[moduli-section-phantom]{Champs de modules}
\item \hyperref[moduli-curves-section-phantom]{Espaces de modules de courbes}
\end{enumerate}
Divers
\begin{enumerate}
\setcounter{enumi}{109}
\item \hyperref[examples-section-phantom]{Exemples}
\item \hyperref[exercises-section-phantom]{Exercices}
\item \hyperref[guide-section-phantom]{Guide bibliographique}
\item \hyperref[desirables-section-phantom]{Desiderata}
\item \hyperref[coding-section-phantom]{Style de programmation}
\item \hyperref[obsolete-section-phantom]{Obsolète}
\item \hyperref[fdl-section-phantom]{Licence de documentation libre GNU}
\item \hyperref[index-section-phantom]{Index généré automatiquement}
\end{enumerate}
\end{multicols}


\bibliography{my}
\bibliographystyle{amsalpha}

\end{document}
